% Предложение темы магистерской диссертации (русский). Сборка: make proposal → thesis/build/proposal-ru.pdf
% Факты и числа — из репозитория проекта (docs/, src/); источники — SRC-NNN в sources.bib.
\documentclass[a4paper,11pt]{article}
\newcommand{\mainlang}{russian}
\newcommand{\docstyle}{article}
% Преамбула: шрифты (латиница, кириллица, арабский с огласовками), поля, библиография.
\usepackage{etoolbox}
\usepackage{fontspec}
\usepackage{polyglossia}
\setdefaultlanguage{\mainlang}
\ifdefstring{\mainlang}{english}{\setotherlanguages{russian,arabic}}{\setotherlanguages{english,arabic}}

% Times New Roman есть в macOS; на другой машине — Liberation Serif или Tempora.
\setmainfont{Times New Roman}
\newfontfamily\cyrillicfont{Times New Roman}
% Арабский — тот же шрифт, что рисует карточки бота (ADR-0024): огласовки и шадда не теряются.
\newfontfamily\arabicfont[Script=Arabic,Path=../src/muallim/presentation/telegram/assets/fonts/,
  Extension=.ttf,UprightFont=*-Regular,BoldFont=*-Bold,Scale=1.25]{ScheherazadeNew}
% \ar{…} — арабская вставка в строке; источник строки — в тексте рядом ([SRC-NNN]).
\newcommand{\ar}[1]{\textarabic{#1}}

\usepackage[left=3cm,right=1.5cm,top=2cm,bottom=2cm]{geometry} % [утвердить] по требованиям вуза
\usepackage{setspace}
\onehalfspacing
\usepackage{indentfirst}

\usepackage{graphicx}
\graphicspath{{../docs/architecture/diagrams/}{../docs/design/screens/}{figures/}}
\usepackage{booktabs,tabularx,longtable}
\usepackage{caption}
\usepackage{csquotes}
\usepackage[hidelinks]{hyperref}

% Библиография — тот же sources.bib, что у всего проекта (ключи SRC-NNN).
\ifdefstring{\mainlang}{english}
  {\usepackage[backend=biber,style=ieee,sorting=none]{biblatex}}
  {\usepackage[backend=biber,style=gost-numeric,sorting=none]{biblatex}}
\addbibresource{../docs/research/sources.bib}

% ID рисунков и таблиц из docs/thesis/figures.md: \label{fig:R-01}, \label{tab:T-01}.

% Данные для титульных листов всех документов. [утвердить] — с руководителем и учебным офисом.
\newcommand{\thesisProgramCode}{09.04.01}                     % [утвердить] код направления
\newcommand{\thesisProgramRu}{Информатика и вычислительная техника} % [утвердить]
\newcommand{\thesisProgramEn}{Computer Science}                % [утвердить]
\newcommand{\thesisSpecialtyRu}{Программная инженерия}
\newcommand{\thesisSpecialtyEn}{Software Engineering}
% Тема — предложение автора (2026-10-18; с ADR-0051 — с адаптацией модели), утверждает руководитель [утвердить].
\newcommand{\thesisTopicRu}{ИИ-тренажёр чтения по букварю «Муаллим ас-Сани»: адаптация открытой модели распознавания фонем к речи русскоязычных начинающих и оценка обратной связи}
\newcommand{\thesisTopicEn}{An AI Reading Tutor for the \emph{Muallim al-Sani} Primer: Adapting an Open Phoneme Recognition Model to Russian-Speaking Beginners and Evaluating Pronunciation Feedback}
\newcommand{\thesisShortRu}{ИИ-тренажёр «Муаллим ас-Сани» и адаптация модели}
\newcommand{\thesisShortEn}{AI Tutor for \emph{Muallim al-Sani} and Model Adaptation}
\newcommand{\thesisAuthorRu}{Исмаилов Абдирахим Ерланович}
\newcommand{\thesisAuthorEn}{Abdirakhim Ismailov}               % латиница — [утвердить] как в паспорте
\newcommand{\thesisAuthorEnFull}{Ismailov Abdirakhim Erlanovich}
\newcommand{\thesisSupervisorRu}{Аскарбекулы Нурсултан}
\newcommand{\thesisSupervisorEn}{Nursultan Askarbekuly}
\newcommand{\thesisUniversityRu}{Университет Иннополис}
\newcommand{\thesisUniversityEn}{Innopolis University}
\newcommand{\thesisCity}{Иннополис}
\newcommand{\thesisCityEn}{Innopolis}
\newcommand{\thesisYear}{2027}

\hypersetup{pdftitle={Предложение темы: \thesisShortRu},pdfauthor={\thesisAuthorRu}}
\fancyhead[R]{\small\sffamily\color{inkgray}Предложение темы · А.\,Е.~Исмаилов}
\setcounter{secnumdepth}{2}

\begin{document}
\thispagestyle{plain}
\begin{center}
{\sffamily\small\color{inkgray}\thesisUniversityRu{} · магистратура «\thesisSpecialtyRu» · предложение темы ВКР}\\[10pt]
{\sffamily\LARGE\bfseries\color{accent}\thesisTopicRu\par}
\vspace{10pt}
{\color{gold}\rule{0.35\linewidth}{1pt}}\\[8pt]
{\large \thesisAuthorRu}\\[2pt]
{\small Научный руководитель: \thesisSupervisorRu}\\[2pt]
{\small\color{inkgray} Черновик для обсуждения с руководителем · октябрь 2026 · защита — июнь 2027}
\end{center}

\begin{keybox}[Кратко]
Чтение Корана начинается с букваря и учителя, который слушает каждый слог. В работе создаётся
Telegram-тренажёр, который ведёт ученика по классическому букварю Ахмад-Хади Максуди «Муаллим ас-Сани»
страница за страницей, слушает голосовое сообщение, находит ошибки произношения на уровне фонем и
объясняет каждую ошибку по-русски. Тренажёр одновременно собирает с согласия учеников первый набор
букварного чтения русскоязычных начинающих с разметкой преподавателей. На этом наборе открытая модель
Quran Muaalem \emph{адаптируется} к новой аудитории двумя способами: калиброванным верификатором с
исходом «не уверен» и дообучением низкого ранга (LoRA). Работа отвечает на четыре вопроса: насколько
готовая модель согласуется с преподавателями и какие её «ошибки ученика» на самом деле ошибки модели
(RQ-1); насколько адаптация на собственном наборе закрывает разрыв с учителем на учениках, которых
модель не видела, не ухудшая чтение экспертов (RQ-2, центральный); помогают ли объяснения следующей
попытке и запоминанию (RQ-3); что предотвращает и чего стоит обвязка ИИ-агентов, в которой создавалась
и обучалась система (RQ-4). Ожидаемый вклад --- измерение и адаптация открытой модели проверки чтения
Корана для взрослых, для которых арабский не родной, с замороженным тестом и слепой разметкой; набор
данных и воспроизводимый конвейер обучения как основа докторского исследования; доказательное описание
разработки исследовательского ПО ИИ-агентами под механическим контролем.
\end{keybox}

\section{Актуальность}

Миллионы взрослых в России и Средней Азии хотят читать Коран по-арабски, но выросли с кириллицей
и без звуков, которые в арабском различают смысл. Традиционный путь --- букварь и учитель. Учитель
произносит слог, ученик повторяет, учитель сразу поправляет каждую оговорку. Такая устная передача,
\emph{таляккы}, работает хорошо, но плохо масштабируется: учитель слушает одного ученика за раз, а
многие ученики живут далеко от медресе или занимаются после работы.

В этом регионе больше века господствует один букварь. Ахмад-Хади Максуди (1868--1941) написал
«Муаллим ас-Сани» («Второй учитель») как мост от знания букв к правильному чтению Корана. Книга
построена на звуковом методе джадидов: буква изучается как звук, а не как название, и слоги
постепенно складываются в слова \cite{SRC-015,SRC-016,SRC-017}. Её до сих пор переиздают, а
духовные управления мусульман Татарстана и Казахстана официально её распространяют
\cite{SRC-020,SRC-040}.

Тем временем речевые технологии научились слушать на уровне фонем. Открытая модель Quran Muaalem
(2025) распознаёт фонемы вместе с десятью артикуляционными признаками (\emph{сифат}) и на чтении
экспертов ошибается в 0,54\,\% фонем \cite{SRC-002}. Метаанализы показывают, что автоматическая
обратная связь по произношению даёт средний положительный эффект ($g=0{,}69$), сильнее всего --- на
отдельных звуках и при явном указании ошибки \cite{SRC-025}. Это ровно задача букваря. Возможность
очевидна: дать каждому начинающему терпеливого слушателя между занятиями с живым учителем.

Препятствие так же очевидно. Muaalem обучена только на чтении мужчин-экспертов. На наборе
qdat\_bench с настоящими ошибками доля ошибок по фонемам вырастает до 5,78\,\%, а F1 для правила
\emph{ихфа} падает до 0,45 \cite{SRC-002}. Лучшая система конкурса IQRA 2026 достигает F1 0,72 по
фонемам, и организаторы отмечают, что в данных нет изучающих арабский как второй язык
\cite{SRC-031,SRC-086}. На речи с акцентом автоматические системы дают ложные замечания, поэтому
обзоры советуют оставлять учителя в контуре \cite{SRC-026}. Как такая модель ведёт себя на
русскоязычном взрослом, который впервые медленно читает отдельные слоги, никто не измерял, и
никто не пробовал её к такому ученику приспособить. Между тем приспособить модель к новой аудитории
можно и на малых данных: дообучение самообучаемых речевых энкодеров на речи изучающих язык улучшает
обнаружение ошибок \cite{SRC-116}, дообучение низкого ранга меняет лишь малую долю весов
\cite{SRC-111}, а калибровка и отказ от ответа позволяют модели говорить «не уверен» вместо ложного
вердикта \cite{SRC-112,SRC-117}.

\section{Пробел в исследованиях}

Обзор литературы, приложений и наборов данных, выполненный для этого предложения \cite{SRC-038},
выявил пять пробелов, которые можно закрыть в сроки магистратуры.

\begin{enumerate}
  \item \textbf{Нет оценки на начинающих букварного уровня.} Наборы для оценки чтения Корана
    (QuranMB, Iqra'Eval, IQRA 2026) содержат беглое чтение аятов \cite{SRC-004,SRC-043,SRC-031}.
    В ближайшем наборе чтения неарабов, RetaSy, нет носителей русского и казахского \cite{SRC-032}.
  \item \textbf{Нет исследований ошибок русско- и казахоязычных учеников.} Типичные трудности
    (фарингальные, эмфатические, межзубные, долгота гласных, удвоение) описаны в методических
    источниках и в работах о других родных языках \cite{SRC-023,SRC-034,SRC-035}, но для этой
    аудитории не измерены.
  \item \textbf{Нет адаптации моделей к L2-начинающим.} Открытые модели чтения Корана обучены на
    экспертах \cite{SRC-002}; работ, где модель дообучают на чтении неарабов-новичков и честно меряют
    выигрыш на новых учениках, не найдено, а в ML-исследованиях с малыми данными часты утечки между
    обучением и проверкой \cite{SRC-122}.
  \item \textbf{Нет цифровой адаптации букваря с автоматической проверкой} и вообще нет
    контролируемых испытаний букварных методик чтения Корана \cite{SRC-038}. Telegram в обучении
    произношению изучали только как канал, а не как автоматического наставника \cite{SRC-092}.
  \item \textbf{Мало данных о разработке ПО ИИ-агентами.} Известно, что агенты подгоняют решение под
    собственные тесты \cite{SRC-027,SRC-082}, а файлы контекста репозитория почти не влияют на
    корректность \cite{SRC-074,SRC-077}. Практики описывают «обвязку» агента (harness engineering)
    \cite{SRC-029,SRC-028}, но работ о том, как она ведёт себя на протяжении целого
    исследовательского проекта, почти нет \cite{SRC-073}.
\end{enumerate}

\section{Цель, задачи и исследовательские вопросы}

\textbf{Цель.} Создать тренажёр чтения, точно следующий «Муаллим ас-Сани», собрать с его помощью
размеченный набор чтения русскоязычных начинающих и адаптировать на этом наборе открытую модель
распознавания фонем так, чтобы она надёжнее выполняла роль слушающего учителя, --- измерив, какую часть
разрыва с преподавателем закрывает адаптация и где разрыв остаётся.

\textbf{Задачи.}
\begin{enumerate}[label=З\arabic*.]
  \item Перевести букварь в цифровую программу в порядке самой книги так, чтобы каждое упражнение
    прослеживалось до страницы и строки конкретного издания \cite{SRC-040}.
  \item Создать тренажёр, который оценивает чтение на уровне фонем, объясняет ошибки по-русски для
    ученика «с нуля» и сохраняет полный след решения по каждой попытке.
  \item Собрать набор данных: согласие (в том числе на обучение модели) и псевдонимы, пилот с 20--40
    учениками в течение 4--6 недель, разметку по фонемам двумя преподавателями, слепую разметку
    тестовых учеников, карточку набора \cite{SRC-123}.
  \item Построить воспроизводимый конвейер обучения и оценки: разбиение по ученикам с замороженным
    тестом, калиброванный верификатор B1, дообучение LoRA B2 с пополнением данных по минимальным парам
    книги, защищённый оценщик.
  \item Провести внутрисубъектный эксперимент с объяснениями и отложенный тест.
  \item Ответить на RQ-1--RQ-3 с размерами эффекта и доверительными интервалами по
    предрегистрированному плану.
  \item Описать обвязку ИИ-агента, в которой создана и обучена система, и оценить её как кейс-стади (RQ-4).
\end{enumerate}

\begin{rqbox}
\textbf{RQ-1 (исходная модель).} Насколько готовая проверка B0 --- модель Muaalem, выравнивание по
правилам и пороги --- согласуется с преподавателями на чтении начинающих по «Муаллим ас-Сани» и какие из
найденных ею ошибок принадлежат ученику, а не модели?\\[4pt]
\textbf{RQ-2 (адаптация, центральный).} Насколько адаптация модели на собранном наборе ---
калиброванный верификатор с исходом «не уверен» (B1) и дообучение низкого ранга (B2) --- уменьшает
ложные отклонения и приближает проверку к преподавателям на учениках, не участвовавших в обучении,
не ухудшая распознавание чтения экспертов?\\[4pt]
\textbf{RQ-3 (обратная связь).} Улучшают ли явные объяснения на русском следующую попытку и
отложенное запоминание по сравнению с простым сигналом «верно / неверно»?\\[4pt]
\textbf{RQ-4 (методика).} Какие известные сбои разработки и ML-экспериментов с ИИ-агентами
предотвращает или делает видимыми механически обеспеченная обвязка и сколько вмешательств человека она
требует?
\end{rqbox}

Гипотезы ниже записаны в предрегистрации, которая будет заморожена отпечатком SHA-256 до начала
сбора данных \cite{SRC-067}. Всё остальное подаётся как разведочный анализ.

\begin{table}[ht]
\centering\small
\caption{Подтверждающие гипотезы (предрегистрация v0.3, утверждается с руководителем).}
\label{tab:hyp-ru}
\begin{tabularx}{\linewidth}{@{}lL@{}}
\toprule
ID & Гипотеза \\
\midrule
H1.1 & Precision и recall обнаружения ошибок B0 на начинающих ниже, чем на домене обучения модели
       (чтение Корана). \\
H1.2 & У B0 ложные «не засчитано» (FRR) встречаются чаще ложных «засчитано» (FAR). \\
H1.3 & Согласие «B0 --- учитель» ($\kappa$) ниже согласия «учитель --- учитель». \\
H1.4 & У русскоязычных больше подтверждённых ошибок в твёрдости (\emph{тафхим}), чем в месте
       образования звука. \\
H2.1 & На тестовых учениках B1 даёт меньший FRR, чем B0, при recall не ниже, чем у B0, минус 0,05. \\
H2.2 & На тестовых учениках F1 обнаружения ошибок по фонемам у B2 выше, чем у B0. \\
H2.3 & B2 не хуже B0 на чтении экспертов: доля ошибок по фонемам растёт не больше чем на 1 п.\,п. \\
H3.1 & После полного объяснения прирост балла на следующей попытке больше, чем после сигнала
       «верно / неверно». \\
H3.2 & Разница H3.1 сохраняется на отложенном тесте через 1--2 недели. \\
\bottomrule
\end{tabularx}
\end{table}

\paragraph{Границы работы} Программа охватывает всю книгу (с.~10--85), но оценка сосредоточена на
уроках букв и огласовок (с.~10--24): с них начинают и в них больше всего ошибок. Правила таджвида
вроде \emph{ихфа} и \emph{идгам}, другие родные языки, активное обучение и персонализация модели
под ученика оставлены для докторантуры. Модель адаптируется, а не обучается с нуля: данных пилота на
порядки меньше, чем нужно для обучения с нуля. Тренажёр дополняет учителя и не претендует его заменить.

\section{Подход к исследованию}

Работа выполнена в парадигме design science research \cite{SRC-100,SRC-101}: создаются артефакты,
решающие реальную проблему, и каждый оценивается методом, подходящим к конкретному утверждению.
Артефактов три (рис.~\ref{fig:design-ru}). Артефакт~A --- тренажёр чтения вместе с набором данных;
он оценивается измерительным исследованием и пилотным экспериментом. Артефакт~C --- адаптированная
модель и конвейер её обучения; он оценивается контролируемым экспериментом на замороженном тесте.
Артефакт~B --- обвязка разработки; он оценивается как встроенное кейс-стади по Рунесону и Хёсту
\cite{SRC-102}.
Для самопроверки каждого типа исследования используются эмпирические стандарты ACM SIGSOFT
\cite{SRC-103}.

\begin{figure}[ht]
\centering
\resizebox{\linewidth}{!}{% Дизайн исследования: три артефакта, четыре RQ, данные и тип исследования.
% Источник: docs/research/metrics-gqm.md (G1–G7), preregistration.md (v0.3), roadmap 1.4–1.8, ADR-0051.
\begin{tikzpicture}[node distance=4mm and 6mm]
  \node[hot,text width=4.1cm] (a1) {\textbf{\LANG{Artefact A}{Артефакт A}}: \LANG{reading tutor}{тренажёр чтения}\\[-1pt]{\scriptsize\LANG{book-ordered curriculum, B0 assessment, feedback, dataset}{программа по книге, оценка B0, объяснения, набор данных}}};
  \node[hot,right=5mm of a1,text width=4.1cm] (a3) {\textbf{\LANG{Artefact C}{Артефакт C}}: \LANG{adapted model}{адаптированная модель}\\[-1pt]{\scriptsize\LANG{B1 calibrated verifier, B2 LoRA fine-tuning}{B1 калиброванный верификатор, B2 дообучение LoRA}}};
  \node[goldbox,right=5mm of a3,text width=4.1cm] (a2) {\textbf{\LANG{Artefact B}{Артефакт B}}: \LANG{agent harness}{обвязка агента}\\[-1pt]{\scriptsize\LANG{hooks, frozen tests and test set, reviewers, journals}{хуки, замороженные тесты и тест, ревьюеры, журналы}}};

  \node[rqbox-node,minimum height=1.25cm,below=7mm of a1,xshift=-2.0cm,text width=2.9cm] (rq1) {\textbf{RQ-1}\\[-1pt]{\scriptsize\LANG{baseline vs.\ teachers; error profile}{B0 против учителей; профиль ошибок}}};
  \node[rqbox-node,minimum height=1.25cm,right=of rq1,text width=2.9cm] (rq2) {\textbf{RQ-2}\\[-1pt]{\scriptsize\LANG{does adaptation close the gap?}{закрывает ли адаптация разрыв?}}};
  \node[rqbox-node,minimum height=1.25cm,right=of rq2,text width=2.9cm] (rq3) {\textbf{RQ-3}\\[-1pt]{\scriptsize\LANG{effect of explanations}{эффект объяснений}}};
  \node[rqbox-node,minimum height=1.25cm,right=of rq3,text width=2.9cm] (rq4) {\textbf{RQ-4}\\[-1pt]{\scriptsize\LANG{what the harness prevents, at what cost}{что предотвращает обвязка и какой ценой}}};

  \node[soft,minimum height=1.75cm,below=of rq1,text width=2.9cm] (d1) {\scriptsize\LANG{phoneme-level annotation; P/R/F1, FAR/FRR, $\kappa$; confirmed confusions}{разметка по фонемам; P/R/F1, FAR/FRR, $\kappa$; подтверждённые путаницы}};
  \node[soft,minimum height=1.75cm,below=of rq2,text width=2.9cm] (d2) {\scriptsize\LANG{learner-disjoint frozen test, blind gold labels; B0 vs.\ B1 vs.\ B2}{замороженный тест по ученикам, слепое «золото»; B0, B1, B2}};
  \node[soft,minimum height=1.75cm,below=of rq3,text width=2.9cm] (d3) {\scriptsize\LANG{within-subject: full vs.\ minimal feedback; delayed test}{внутрисубъектно: полное vs.\ краткое; отложенный тест}};
  \node[soft,minimum height=1.75cm,below=of rq4,text width=2.9cm] (d4) {\scriptsize\LANG{embedded case study: hook logs, sessions, commits}{встроенное кейс-стади: журналы хуков, сессии, коммиты}};

  \node[lab,below=1mm of d1] {\LANG{measurement study}{измерительное}};
  \node[lab,below=1mm of d2] {\LANG{controlled ML experiment}{контролируемый ML-эксперимент}};
  \node[lab,below=1mm of d3] {\LANG{pilot experiment}{пилотный эксперимент}};
  \node[lab,below=1mm of d4] {\LANG{case study}{кейс-стади}};

  \draw[arr] (a1.south) -- (rq1.north);
  \draw[arr] (a1.south) -- (rq3.north);
  \draw[arr] (a1.south) -- (rq2.north);
  \draw[arr] (a3.south) -- (rq2.north);
  \draw[arr] (a2.south) -- (rq4.north);
  \foreach \r/\d in {rq1/d1,rq2/d2,rq3/d3,rq4/d4}{\draw[arr] (\r) -- (\d);}
\end{tikzpicture}
}
\caption{Дизайн исследования: три артефакта, четыре вопроса и данные для каждого.}
\label{fig:design-ru}
\end{figure}

\subsection{Артефакт A: тренажёр чтения}

Тренажёр уже работает (версия 0.6.0). В программе вся книга --- 125 уроков и 800 карточек в девяти
частях --- в порядке книги: от хамзы и группы букв \ar{ر ز م} на с.~10 до коротких сур в конце. У
каждой карточки указаны страница и строка альметьевского издания 2010~г., а 139 примеров из Корана
сверены с текстом Мединского мусхафа \cite{SRC-040,SRC-041}.

Ученик открывает карточку, может послушать эталон и присылает голосовое. Конвейер оценки
(рис.~\ref{fig:pipe-ru}) переводит запись в нужный формат, распознаёт фонемы моделью Muaalem,
выравнивает их с каноническими фонемами строки для чтения с остановкой и слитно, относит каждое
расхождение к одному из десяти типов ошибок с указанием слога и применяет пороги оценки. Ученик видит
строку с подсвеченной ошибкой, короткое сравнение «как надо $\to$ как у тебя», подсказку по
артикуляции на русском и следующий шаг («послушай, скажи слог медленно, прочитай строку снова»).
Такой разбор опирается на данные о том, что явная и сфокусированная обратная связь работает лучше
простого повтора правильного ответа и длинных списков \cite{SRC-053,SRC-054,SRC-055,SRC-056}.

\begin{figure}[ht]
\centering
\resizebox{\linewidth}{!}{% Конвейер оценки одной попытки: от голосового до объяснения и записи для исследования.
% Код: application/use_cases/recitation.py, infrastructure/asr/muaalem.py, domain/assessment.py,
% domain/policy.py, application/feedback.py (ADR-0007, ADR-0014, ADR-0019).
\begin{tikzpicture}[node distance=5mm and 7mm]
  \node[goldbox,text width=2.7cm] (voice) {\LANG{Voice message}{Голосовое}\\[-1pt]{\scriptsize OGG/Opus, $\le$60\,s}};
  \node[box,right=of voice,text width=2.6cm] (ff) {ffmpeg\\[-1pt]{\scriptsize\LANG{16 kHz mono}{16 кГц, моно}}};
  \node[hot,right=of ff,text width=3.0cm] (asr) {Quran Muaalem\\[-1pt]{\scriptsize\LANG{phonemes + 10 \emph{sifat}}{фонемы + 10 сифат}}};

  \node[goldbox,below=11mm of voice,text width=2.7cm] (card) {\LANG{Book line}{Строка книги}\\[-1pt]{\scriptsize\LANG{page + line, SRC-040}{страница + строка, SRC-040}}};
  \node[box,right=of card,text width=2.6cm] (phon) {\LANG{Phonetiser}{Фонетизатор}\\[-1pt]{\scriptsize\LANG{\emph{wasl} and \emph{waqf}}{васл и вакф}}};
  \node[box,right=of phon,text width=3.0cm] (align) {\LANG{Alignment}{Выравнивание}\\[-1pt]{\scriptsize\LANG{per reference variant}{с каждым эталоном}}};

  \node[box,right=of asr,text width=2.9cm,yshift=-8mm] (err) {\LANG{Typed errors}{Типизированные ошибки}\\[-1pt]{\scriptsize\LANG{10 kinds, syllable}{10 типов, слог}}};
  \node[box,below=5mm of err,text width=2.9cm] (pol) {\LANG{Grading policy}{Политика оценки}\\[-1pt]{\scriptsize\LANG{score, verdict}{балл, вердикт}}};
  \node[human,below=5mm of pol,text width=2.9cm] (fb) {\LANG{Feedback in Russian}{Объяснение по-русски}\\[-1pt]{\scriptsize\LANG{what / how / next}{что / как / что дальше}}};

  \node[store,below=9mm of align,text width=3.1cm,minimum height=1.2cm] (db) {\LANG{Attempt record}{Запись попытки}\\[-1pt]{\scriptsize\LANG{full decision trace}{полный след решения}}};

  \draw[arr] (voice) -- (ff);
  \draw[arr] (ff) -- (asr);
  \draw[arr] (card) -- (phon);
  \draw[arr] (phon) -- (align);
  \draw[arr] (asr.south) -- node[lab,right]{\LANG{heard}{услышано}} (align.north);
  \draw[arr] (align.east) -- ++(4mm,0) |- (err.west);
  \draw[arr] (err) -- (pol);
  \draw[arr] (pol) -- (fb);
  \draw[darr] (align) -- (db);
  \node[lab,below=1mm of db] {\LANG{audio kept only with consent (SHA-256)}{аудио — только с согласия (SHA-256)}};
\end{tikzpicture}
}
\caption{Оценка одной попытки: от голосового до объяснения и записи для исследования.}
\label{fig:pipe-ru}
\end{figure}

В каждой попытке сохраняются варианты эталона, услышанные фонемы, сырой выход модели, пороги,
версии кода, программы и модели, поэтому любой вердикт можно пересчитать позже с другими порогами
или другой моделью. Аудио хранится только с явного отзываемого согласия; ученики обозначены
псевдонимами; удалить все свои данные можно одной кнопкой.

\subsection{Артефакт C: адаптированная модель}

Сравниваются три системы на одних и тех же записях. \textbf{B0} --- Muaalem как есть с правилами и
порогами; так тренажёр работает сейчас. \textbf{B1} --- обученный верификатор поверх выходов B0: по
каждой фонеме он получает уверенность модели, вероятности артикуляционных признаков, меру
\emph{Goodness of Pronunciation} \cite{SRC-115}, длительность и тип материала и оценивает вероятность
настоящей ошибки. Вероятности калибруются \cite{SRC-112}, а порог выбирается конформным контролем риска
так, чтобы доля ложных отклонений не превышала заданную $\varepsilon$ \cite{SRC-113,SRC-114}; попытки,
где уверенности не хватает, получают исход «не уверен» \cite{SRC-117} --- ученику предлагается
перечитать, а попытка уходит учителю. Верификатору хватает сотен примеров, и он обучается на обычном
процессоре. \textbf{B2} --- Muaalem, дообученная методом LoRA \cite{SRC-111} на фонемах, исправленных
преподавателями; так уже дообучали самообучаемые речевые энкодеры для обнаружения ошибок у изучающих
язык \cite{SRC-116}. Чтобы ошибок в обучающих данных было больше, используется сама книга: верное
чтение карточки с \ar{ك} --- готовый пример ошибки для парной карточки с \ar{ق}, и так для каждой
минимальной пары; это вариант пополнения данных заменой фонем \cite{SRC-119}, дополненный изменением
скорости записи \cite{SRC-120}. Дообученная модель проверяется ещё и на чтении экспертов, чтобы
выигрыш на новичках не был куплен забыванием \cite{SRC-118}.

В конце пилота B1 и B2 работают в тренажёре \emph{в тени}: ученик по-прежнему видит вердикт B0, а
вердикты адаптированных систем записываются рядом. Так эксперимент RQ-3 не смешивается со сменой
модели, а новые попытки дают дополнительную проверку после обучения.

\subsection{Артефакт B: обвязка разработки}

Систему создаёт один автор вместе с ИИ-агентом (Claude Code). Важные правила обеспечивает среда, а не
добрая воля агента \cite{SRC-029,SRC-076}. Хуки запрещают править тесты в фазе реализации, не дают
читать персональные данные, требуют подтверждения человека в его зонах (учебный контент, тексты
ученику, пороги, оценщик RQ-1), проверяют, что новая библиотека действительно есть на PyPI, и не дают
агенту закончить работу, пока не проходят архитектурные контракты, проверка типов и тесты.
Субагенты-ревьюеры без права записи проверяют код, архитектуру и соответствие изменения просьбе автора.
Те же механизмы защищают эксперименты обучения: список тестовых учеников и оценщик заморожены
отпечатком SHA-256, а читать метки теста агенту запрещает хук --- по правилам, сформулированным группой
научного руководителя для исследований с агентами \cite{SRC-067}.
Каждая блокировка и каждый запрос подтверждения записываются в журнал. Этот журнал вместе с журналами
сессий и запросов --- данные для RQ-4.

\subsection{Исследования}

\paragraph{Исследование 1 --- согласие (RQ-1)} Два преподавателя независимо размечают каждую
попытку пилота по фонемам: видят ответ тренажёра и исправляют его, где он ошибся. По протоколам
Iqra'Eval и L2-ARCTIC \cite{SRC-063,SRC-064} каждая фонема эталона относится к верному принятию,
верному отклонению, ложному принятию или ложному отклонению; отсюда precision, recall, F1 и точность
диагноза. На уровне вердикта считаются FAR и FRR, а также $\kappa$ Коэна \cite{SRC-104} между
тренажёром и каждым преподавателем в сравнении с $\kappa$ между преподавателями как потолком.
Согласие разметчиков по фонемам измеряется $\alpha$ Криппендорфа \cite{SRC-065}. Доверительные
интервалы строятся бутстрэпом по ученикам, потому что попытки одного ученика зависимы.
Самоисправления и непригодные записи считаются отдельно.

Профиль ошибок входит в то же исследование. Матрицы путаниц «ожидалось $\to$ услышано» строятся
дважды: по всем ошибкам, найденным тренажёром, и только по ошибкам, подтверждённым разметчиками. Разница
показывает, какая часть «профиля ученика» --- шум модели. Минимальные пары самой книги (\ar{ق}/\ar{ك},
\ar{ص}/\ar{س}, \ar{ط}/\ar{ت}, \ar{ح}/\ar{خ}, \ar{ع}/\ar{غ} и др.) служат заранее заданной линзой.
Короткая анкета при входе фиксирует родной язык и опыт чтения как ковариаты.

\paragraph{Исследование 2 --- адаптация модели (RQ-2)} До первого обучения около 30\,\% учеников,
стратифицированно по родному языку и активности, уходят в тест; их список фиксируется отпечатком.
Ни одна запись тестового ученика не участвует ни в обучении, ни в калибровке, ни в выборе порогов:
разбиение по попыткам дало бы утечку голоса диктора \cite{SRC-122}. Попытки тестовых учеников оба
преподавателя размечают \emph{вслепую}, не видя ответа тренажёра, и согласуют расхождения; так
получается «золотая» разметка, которая не тянет сравнение в пользу B0. Гиперпараметры выбираются
только кросс-валидацией по ученикам внутри обучающей части, иначе оценка была бы завышена
\cite{SRC-121}. Тест открывается один раз, когда все системы обучены и заморожены. Сравнение B0, B1 и B2
--- по тем же метрикам, что в исследовании~1, парным кластерным бутстрэпом по ученикам. Сводная мера ---
\emph{доля закрытого разрыва}: на сколько $\kappa$ системы с «золотом» поднялась от B0 к потолку
$\kappa$ «учитель --- учитель». Дополнительно сообщаются ошибка калибровки, кривая «риск --- покрытие»
верификатора и, как разведка, кривая обучения --- как растёт качество B2 с числом обучающих учеников.
Эта кривая прямо отвечает на вопрос докторантуры: сколько данных нужно новой аудитории.

\paragraph{Исследование 3 --- эксперимент с обратной связью (RQ-3)} Для двух независимых групп
20--40 учеников мало (для среднего эффекта нужно около 45 человек в группе), поэтому дизайн
внутрисубъектный, по образцу Neri et al.\ \cite{SRC-054}. Каждый ученик получает полное объяснение на
случайной половине карточек и только сигнал «верно / неверно» на другой половине. Основной исход ---
изменение балла от неудачной попытки к следующей по той же карточке; дополнительный --- балл на
отложенном тесте через 1--2 недели. Эффект оценивается смешанной моделью со случайными эффектами
ученика и карточки \cite{SRC-106}. Для девяти подтверждающих гипотез применяется поправка Холма.

\paragraph{Исследование 4 --- кейс-стади обвязки (RQ-4)} Источники данных --- журнал хуков
(блокировки и запросы подтверждения по причинам), сводки сессий, журнал запросов, коммиты и записи
архитектурных решений. Метрики --- блокировки по типам, вмешательства человека за сессию, итерации до
зелёных проверок, дефекты, которые проверки пропустили, попытки агента обратиться к тестовой выборке. Самооценка продуктивности не используется: она
расходится с измеренной \cite{SRC-030}.

\section{Ожидаемый вклад}

\begin{enumerate}
  \item \textbf{Эмпирический:} первое измерение открытой модели распознавания фонем Корана на
    букварном чтении начинающих, для которых арабский не родной, относительно оценки преподавателей:
    где она ошибается и как часто отклоняет верное чтение.
  \item \textbf{Модель и метод:} адаптация открытой модели к новой аудитории на малых данных ---
    калиброванный верификатор с гарантией на долю ложных отклонений и дообучение LoRA с пополнением по
    минимальным парам книги --- с честной оценкой на новых учениках и проверкой забывания.
  \item \textbf{Данные:} набор чтения начинающих с разметкой по фонемам, слепым «золотым» тестом и
    карточкой набора, собранный с согласия; публикуется вместе с дообученными весами, если позволят
    этика и согласия \cite{SRC-097,SRC-123}.
  \item \textbf{Методика программной инженерии:} доказательное описание обвязки ИИ-агента для
    исследовательского ПО и ML-экспериментов с измеренной пользой, ценой и ограничениями.
  \item \textbf{Артефакт:} тренажёр с открытой архитектурой, где модель, базу данных и мессенджер
    можно заменить, не трогая предметную логику, --- готовая площадка для сравнения моделей.
\end{enumerate}

\paragraph{Задел для докторантуры} Магистерская оставляет докторской работе не только вопросы, но и
инструменты: замороженный тест, который остаётся тестом и дальше, конвейер обучения, кривую обучения и
теневой режим. Отсюда главная докторская линия --- адаптация с человеком в контуре: какие попытки
отдавать учителю на разметку (активное обучение), как использовать жалобы «бот ошибся» как слабые метки,
как персонализировать модель под ученика и переносить её на другие родные языки.

\section{Этика, персональные данные и религиозная чувствительность}

Голосовые записи вместе с религиозной темой делают данные чувствительными. Поэтому согласие на
работу сервиса отделено от согласия на запись голоса; аудио без согласия не хранится; во всех
отчётах и выгрузках --- только псевдонимы; участвовать в исследовании могут только взрослые; удалить
свои данные можно одной кнопкой \cite{SRC-033,SRC-036,SRC-037}. Открытые вопросы
институциональные: кто оператор персональных данных, нужно ли заключение этической комиссии
университета, где расположен сервер. Они будут решены с руководителем до того, как в пилот будут
приглашены люди вне закрытого круга. Тренажёр использует только арабский текст мусхафа и собственные
нейтральные пояснения и не воспроизводит русский текст книги, защищённый авторским правом.
Согласие на запись голоса явно разрешает обучение модели и публикацию весов; голосовые отпечатки не
строятся. Модель Muaalem и её основа w2v-BERT~2.0 распространяются по лицензии MIT \cite{SRC-003},
поэтому дообучение и публикация весов допустимы.

\section{План работ}

План показан на рис.~\ref{fig:gantt-ru}. Подготовка (предложение, этика, сверка программы со сканом
книги, эталонные записи, поиск разметчиков) идёт до декабря 2026~г.; в ноябре--декабре конвейер
обучения отлаживается на открытых данных, чтобы в апреле на данных пилота он только запускался. Сбор
данных --- с января по март 2027~г., список тестовых учеников замораживается в феврале. Обучение B1 и B2
и единственный прогон теста --- в апреле; затем анализ и текст, предзащита в мае и защита в июне 2027~г.

\begin{figure}[ht]
\centering
\resizebox{\linewidth}{!}{% План работ: октябрь 2026 — июнь 2027 (roadmap, этап 1). Месяц = 2 слота (половины месяца).
\begin{ganttchart}[
    hgrid,vgrid={*1{draw=black!8}},
    x unit=0.68cm,y unit chart=0.52cm,y unit title=0.55cm,
    title/.append style={fill=accent!10,draw=accent!40},
    title label font=\sffamily\scriptsize,
    bar/.append style={fill=accent!70,draw=accent},
    bar height=0.55,
    group/.append style={fill=gold,draw=gold},
    group height=0.25,
    milestone/.append style={fill=lesson,draw=lesson},
    bar label font=\sffamily\footnotesize,
    group label font=\sffamily\footnotesize\bfseries,
    milestone label font=\sffamily\footnotesize\itshape,
  ]{1}{18}
  \gantttitle{\LANG{Oct}{Окт}}{2}\gantttitle{\LANG{Nov}{Ноя}}{2}\gantttitle{\LANG{Dec}{Дек}}{2}
  \gantttitle{\LANG{Jan}{Янв}}{2}\gantttitle{\LANG{Feb}{Фев}}{2}\gantttitle{\LANG{Mar}{Мар}}{2}
  \gantttitle{\LANG{Apr}{Апр}}{2}\gantttitle{\LANG{May}{Май}}{2}\gantttitle{\LANG{Jun}{Июн}}{2}\\
  \ganttgroup{\LANG{Preparation}{Подготовка}}{1}{6}\\
  \ganttbar{\LANG{Proposal, preregistration freeze}{Предложение, заморозка предрегистрации}}{1}{3}\\
  \ganttbar{\LANG{Ethics, consent, legal debt (T1)}{Этика, согласие, правовой долг (T1)}}{2}{5}\\
  \ganttbar{\LANG{Book proofreading, reference audio}{Сверка книги, эталонные записи}}{1}{5}\\
  \ganttbar{\LANG{Annotators, pilot readiness}{Разметчики, готовность пилота}}{4}{6}\\
  \ganttgroup{\LANG{Data collection}{Сбор данных}}{7}{12}\\
  \ganttbar{\LANG{Pilot, 20--40 learners, 6 weeks}{Пилот, 20--40 учеников, 6 недель}}{7}{9}\\
  \ganttbar{\LANG{Delayed post-test (RQ-3)}{Отложенный тест (RQ-3)}}{10}{10}\\
  \ganttbar{\LANG{Double annotation (RQ-1)}{Двойная разметка (RQ-1)}}{8}{12}\\
  \ganttbar{\LANG{Agent-harness telemetry (RQ-4)}{Журналы обвязки агента (RQ-4)}}{1}{12}\\
  \ganttgroup{\LANG{Analysis and writing}{Анализ и текст}}{11}{18}\\
  \ganttbar{\LANG{Analysis RQ-1--RQ-4}{Анализ RQ-1--RQ-4}}{11}{14}\\
  \ganttbar{\LANG{Chapters 4--6, abstract}{Главы 4--6, аннотация}}{12}{16}\\
  \ganttmilestone{\LANG{Pre-defence}{Предзащита}}{16}\\
  \ganttbar{\LANG{Revision, plagiarism check}{Доработка, проверка заимствований}}{16}{17}\\
  \ganttmilestone{\LANG{Defence}{Защита}}{18}
\end{ganttchart}
}
\caption{План работ, октябрь 2026 --- июнь 2027 (шаг --- полмесяца).}
\label{fig:gantt-ru}
\end{figure}

\section{Риски}

\begin{table}[ht]
\centering\small
\caption{Главные риски и ответ на них (из реестра рисков проекта).}
\begin{tabularx}{\linewidth}{@{}P{4.3cm}L@{}}
\toprule
Риск & Ответ \\
\midrule
Мало участников пилота & Набирать через преподавателей и мечети заранее; внутрисубъектному дизайну
нужно меньше людей; сообщать размеры эффекта с интервалами, а не только значимость. \\
Модель часто ошибается на начинающих & Честно измеренная низкая точность --- тоже результат RQ-1;
кнопка «Бот ошибся», сырой выход модели в каждой попытке, планируемый исход «не уверен». \\
Нет второго разметчика & Договориться с двумя преподавателями до пилота; запасной вариант ---
автор и один преподаватель, оба размечают все попытки. \\
Правовые и этические препятствия & Обсудить вопросы оператора и этики с руководителем сейчас;
сначала закрыть все пункты техдолга, блокирующие пилот. \\
Мало данных для дообучения & B1 работает на сотнях примеров; пополнение по минимальным парам; если к
15 апреля меньше \textasciitilde1\,500 размеченных попыток, B2 подаётся как проверка осуществимости, а
RQ-2 отвечает B1 (зафиксировано в предрегистрации). \\
Утечка между обучением и тестом & Разбиение по ученикам, список теста и оценщик заморожены отпечатком,
хук запрещает агенту читать метки теста, тест открывается один раз. \\
Нет GPU & Кластер университета или облако; LoRA требует одной карты; B1 обучается на процессоре. \\
Ошибки в арабском тексте программы & Сверка человеком со сканом; версия программы в каждой
попытке позволяет найти и исключить затронутые попытки. \\
Объём больше срока & Приоритет --- RQ-1 и RQ-2 на с.~10--24; приоритеты MoSCoW; пересмотр на каждой
встрече с руководителем. \\
\bottomrule
\end{tabularx}
\end{table}

\section{Текущее состояние и предварительные данные}

Система работает целиком, с настоящей моделью. В первой проверке на обучающем чтении Махмуда Халиля
аль-Хусари \cite{SRC-042} верные аяты получили 99--100 баллов, а намеренно неверно сопоставленные ---
40--67: на чтении эксперта конвейер отличает верное чтение от неверного. На девяти первых попытках
медианное время обработки составило 1,7\,с (90-й процентиль --- 5,0\,с), что укладывается в цель
10\,с, хотя сюда не входят загрузка файла и отрисовка ответа. Свойства алгоритма оценки проверены на
5\,000 случайных строк фонем на каждое свойство --- нарушений нет \cite{SRC-096}. В коде около
9\,000 строк на Python в 92 модулях, 736 автоматических тестов проходят, записано 50 архитектурных
решений, 47 функциональных требований прослежены до тестов. Каждая попытка уже хранит сырой выход
модели с уверенностью по фонемам и версию модели --- ровно то, что нужно для обучения B1 и сравнения
систем. Ни одно из этих чисел пока не отвечает на исследовательский вопрос; они показывают, что
инструмент готов.

\section{Предлагаемая структура диссертации}

\begin{enumerate}[leftmargin=2em]
  \item \textbf{Введение} --- контекст, проблема, цель, вопросы, вклад.
  \item \textbf{Обзор литературы} --- букварь и звуковой метод; компьютерное обучение произношению;
    обнаружение ошибок чтения Корана; адаптация речевых моделей и калибровка; программная инженерия
    с ИИ-агентами.
  \item \textbf{Методология} --- design science; требования и сценарии качества; метрики GQM;
    дизайн исследований; разбиение данных и защита теста; предрегистрация; угрозы валидности.
  \item \textbf{Реализация} --- программа по книге; архитектура; оценка и объяснения; данные и
    разметка; конвейер обучения B1 и B2; приватность; обвязка агента.
  \item \textbf{Оценка и обсуждение} --- RQ-1--RQ-4, ограничения.
  \item \textbf{Заключение} --- ответы, вклад, направления докторантуры.
\end{enumerate}

\printbibliography[heading=bibintoc,title={Список источников}]
\end{document}
